\documentclass[12pt,a4paper]{article}

% ------------------------------------------------------------
% Кодировка и русский язык
% ------------------------------------------------------------
\usepackage[T2A]{fontenc}
\usepackage[utf8]{inputenc}
\usepackage[russian,english]{babel}

% ------------------------------------------------------------
% Математика
% ------------------------------------------------------------
\usepackage{amsmath,amssymb,amsthm,mathtools,bm}
\usepackage{array,booktabs,tabularx,multirow}
\usepackage{enumitem}
\usepackage[expansion=false]{microtype}

% ------------------------------------------------------------
% Оформление
% ------------------------------------------------------------
\usepackage[
    left=2.2cm,
    right=2.2cm,
    top=2.0cm,
    bottom=2.0cm
]{geometry}
\usepackage{xcolor}
\usepackage[most]{tcolorbox}
\usepackage{tikz}
\usetikzlibrary{positioning,arrows.meta,calc}
\usepackage{hyperref}

\hypersetup{
    colorlinks=true,
    linkcolor=blue!45!black,
    urlcolor=blue!55!black,
    citecolor=blue!45!black
}

\setlength{\parindent}{1.25cm}
\setlength{\parskip}{0.2em}
\setlist{nosep}

% ------------------------------------------------------------
% Обозначения
% ------------------------------------------------------------
\newcommand{\R}{\mathbb{R}}
\newcommand{\E}{\mathbb{E}}
\newcommand{\KL}{\operatorname{KL}}
\newcommand{\JS}{\operatorname{JS}}
\newcommand{\softmax}{\operatorname{softmax}}
\newcommand{\ReLU}{\operatorname{ReLU}}
\newcommand{\Lip}{\operatorname{Lip}}
\newcommand{\tr}{\operatorname{tr}}
\newcommand{\diag}{\operatorname{diag}}
\newcommand{\sgn}{\operatorname{sgn}}
\newcommand{\argminop}{\operatorname*{arg\,min}}
\newcommand{\argmaxop}{\operatorname*{arg\,max}}

\numberwithin{equation}{section}

% ------------------------------------------------------------
% Теоремы
% ------------------------------------------------------------
\newtheorem{theorem}{Теорема}[section]
\newtheorem{proposition}[theorem]{Предложение}
\newtheorem{lemma}[theorem]{Лемма}
\newtheorem{corollary}[theorem]{Следствие}

\theoremstyle{definition}
\newtheorem{definition}[theorem]{Определение}
\newtheorem{example}[theorem]{Пример}

\theoremstyle{remark}
\newtheorem{remark}[theorem]{Замечание}

\renewcommand{\proofname}{Доказательство}

% ------------------------------------------------------------
% Цветные блоки
% ------------------------------------------------------------
\newtcolorbox{algobox}[1]{
    enhanced,
    breakable,
    colback=blue!3,
    colframe=blue!45!black,
    title={#1},
    fonttitle=\bfseries,
    boxrule=0.7pt,
    arc=2mm
}

\newtcolorbox{examplebox}[1]{
    enhanced,
    breakable,
    colback=green!3,
    colframe=green!40!black,
    title={#1},
    fonttitle=\bfseries,
    boxrule=0.7pt,
    arc=2mm
}

\newtcolorbox{importantbox}[1]{
    enhanced,
    breakable,
    colback=red!3,
    colframe=red!50!black,
    title={#1},
    fonttitle=\bfseries,
    boxrule=0.7pt,
    arc=2mm
}

\tikzset{
    block/.style={
        draw,
        rounded corners,
        align=center,
        minimum height=10mm,
        minimum width=22mm,
        fill=blue!5,
        font=\small
    },
    arrow/.style={-{Latex[length=2mm]},thick}
}

% ------------------------------------------------------------
% Документ
% ------------------------------------------------------------

\title{Билет №98 \\ \small Сверточные сети. Генеративные состязательные сети. (ИТМО)}
\author{Сериков Владислав Александрович}
\date{19 июля 2026}

\begin{document}

\maketitle
\tableofcontents
\newpage

% ============================================================
\section{Основные обозначения}
% ============================================================

Вектор обозначается строчной полужирной буквой, например
\(\bm{x}\in\R^n\), матрица --- прописной буквой \(W\), а многомерный
тензор --- буквой \(X\).

Для изображения будем использовать представление
\[
    X\in\R^{C_{\mathrm{in}}\times H\times W},
\]
где \(C_{\mathrm{in}}\) --- число входных каналов, \(H\) --- высота,
\(W\) --- ширина. Например, для RGB-изображения
\(C_{\mathrm{in}}=3\).

Обучающая выборка имеет вид
\[
    \mathcal{D}=\{(\bm{x}_i,y_i)\}_{i=1}^{N}
\]
в задаче обучения с учителем или
\[
    \mathcal{D}=\{\bm{x}_i\}_{i=1}^{N}
\]
в задаче обучения без учителя.

Параметры нейронной сети обозначаются через \(\theta\). Обучение обычно
сводится к минимизации эмпирического риска
\[
    \widehat{R}(\theta)
    =
    \frac{1}{N}\sum_{i=1}^{N}
    \ell\bigl(f_\theta(\bm{x}_i),y_i\bigr),
\]
где \(\ell\) --- функция потерь.

% ============================================================
\section{Сверточные нейронные сети}
% ============================================================

\subsection{Идея сверточной сети}

\begin{definition}
\emph{Сверточная нейронная сеть} (convolutional neural network, CNN) ---
нейронная сеть, в которой хотя бы часть линейных преобразований
реализуется сверточными слоями, использующими локальные связи и
разделение параметров между пространственными позициями.
\end{definition}

Основные предпосылки CNN:

\begin{enumerate}
    \item \textbf{Локальность.}
    Значение признака в некоторой позиции зависит от небольшой
    окрестности входа.

    \item \textbf{Разделение параметров.}
    Один и тот же фильтр применяется во всех пространственных позициях.

    \item \textbf{Иерархия признаков.}
    Ранние слои обнаруживают простые локальные структуры, например
    границы и углы, а глубокие слои объединяют их в более сложные
    признаки.

    \item \textbf{Эквивариантность относительно сдвига.}
    Сдвиг входного объекта приводит к соответствующему сдвигу карты
    признаков, если не учитывать границы и уменьшение разрешения.
\end{enumerate}

Типичная схема классификационной CNN показана ниже.

\begin{center}
\begin{tikzpicture}[node distance=7mm]
    \node[block] (input) {Изображение};
    \node[block,right=of input] (conv) {Свертка};
    \node[block,right=of conv] (act) {Нелинейность};
    \node[block,right=of act] (pool) {Уменьшение\\разрешения};
    \node[block,right=of pool] (features) {Карты\\признаков};
    \node[block,right=of features] (head) {Классификатор};

    \draw[arrow] (input) -- (conv);
    \draw[arrow] (conv) -- (act);
    \draw[arrow] (act) -- (pool);
    \draw[arrow] (pool) -- (features);
    \draw[arrow] (features) -- (head);
\end{tikzpicture}
\end{center}

% ------------------------------------------------------------
\subsection{Дискретная свертка и взаимная корреляция}
% ------------------------------------------------------------

Для одномерных последовательностей математическая дискретная свертка
определяется как
\[
    (x*k)[n]
    =
    \sum_{m\in\mathbb{Z}} x[m]k[n-m].
\]

В программных библиотеках глубокого обучения обычно реализуется
\emph{взаимная корреляция}:
\[
    (x\star k)[n]
    =
    \sum_{m\in\mathbb{Z}} k[m]x[n+m].
\]

От свертки она отличается отсутствием отражения ядра. Поскольку ядро
является обучаемым, это различие не влияет на выразительность слоя.
Поэтому взаимную корреляцию в контексте нейронных сетей обычно также
называют сверткой.

Рассмотрим вход
\[
    X\in\R^{C_{\mathrm{in}}\times H\times W}
\]
и ядро
\[
    K\in
    \R^{
        C_{\mathrm{out}}
        \times C_{\mathrm{in}}
        \times k_h\times k_w
    }.
\]

После дополнения входа нулями обозначим его через
\(\widetilde{X}\). Тогда выход сверточного слоя имеет вид
\[
\begin{split}
    Y_{o,i,j}
    &=
    b_o
    +
    \sum_{c=0}^{C_{\mathrm{in}}-1}
    \sum_{u=0}^{k_h-1}
    \sum_{v=0}^{k_w-1}
    K_{o,c,u,v}
    \widetilde{X}_{
        c,\,
        i s_h+u d_h,\,
        j s_w+v d_w
    },
\end{split}
\]
где

\begin{itemize}
    \item \(o\) --- номер выходного канала;
    \item \(s_h,s_w\) --- шаги свертки;
    \item \(d_h,d_w\) --- коэффициенты дилатации;
    \item \(b_o\) --- смещение выходного канала.
\end{itemize}

После линейной операции обычно применяется нелинейность:
\[
    Z_{o,i,j}=\varphi(Y_{o,i,j}).
\]

% ------------------------------------------------------------
\subsection{Размер выхода сверточного слоя}
% ------------------------------------------------------------

\begin{definition}
Эффективный размер ядра с дилатацией определяется формулами
\[
    k_h^{\mathrm{eff}}=d_h(k_h-1)+1,
    \qquad
    k_w^{\mathrm{eff}}=d_w(k_w-1)+1.
\]
\end{definition}

\begin{proposition}[Размер выхода свертки]
Пусть вход имеет пространственный размер \(H\times W\), дополнение
нулями равно \(p_h,p_w\), шаг равен \(s_h,s_w\), а дилатация ---
\(d_h,d_w\). Если дополненный вход не меньше эффективного ядра, то
\[
    H_{\mathrm{out}}
    =
    \left\lfloor
    \frac{
        H+2p_h-d_h(k_h-1)-1
    }{s_h}
    \right\rfloor+1,
\]
\[
    W_{\mathrm{out}}
    =
    \left\lfloor
    \frac{
        W+2p_w-d_w(k_w-1)-1
    }{s_w}
    \right\rfloor+1.
\]
\end{proposition}

\begin{proof}
Рассмотрим вертикальное направление. После дополнения нулями доступная
высота равна
\[
    H'=H+2p_h.
\]
Эффективная высота ядра равна
\[
    k_h^{\mathrm{eff}}=d_h(k_h-1)+1.
\]

Левая граница первого положения ядра имеет координату \(0\), второго
--- \(s_h\), третьего --- \(2s_h\). В общем случае левая граница
положения с номером \(q\) равна \(q s_h\).

Ядро полностью помещается во вход, если
\[
    q s_h+k_h^{\mathrm{eff}}\leq H'.
\]
Следовательно,
\[
    q
    \leq
    \frac{H'-k_h^{\mathrm{eff}}}{s_h}.
\]
Максимальный допустимый целый номер положения равен
\[
    q_{\max}
    =
    \left\lfloor
    \frac{H'-k_h^{\mathrm{eff}}}{s_h}
    \right\rfloor.
\]
Поскольку нумерация начинается с нуля, число положений равно
\(q_{\max}+1\). Поэтому
\[
\begin{split}
    H_{\mathrm{out}}
    &=
    \left\lfloor
    \frac{
        H+2p_h-\bigl(d_h(k_h-1)+1\bigr)
    }{s_h}
    \right\rfloor+1\\
    &=
    \left\lfloor
    \frac{
        H+2p_h-d_h(k_h-1)-1
    }{s_h}
    \right\rfloor+1.
\end{split}
\]
Для ширины доказательство аналогично.
\end{proof}

\paragraph{Основные режимы дополнения.}

\begin{itemize}
    \item \textbf{Valid convolution:} \(p_h=p_w=0\). Размер выхода
    уменьшается.

    \item \textbf{Same convolution:} при шаге \(1\) выбирается такое
    дополнение, чтобы размер выхода совпал с размером входа. Для
    нечетного ядра и единичной дилатации:
    \[
        p_h=\frac{k_h-1}{2},
        \qquad
        p_w=\frac{k_w-1}{2}.
    \]

    \item \textbf{Causal convolution:} в одномерных временных моделях
    используются только текущие и предыдущие значения входа.
\end{itemize}

% ------------------------------------------------------------
\subsection{Численный пример свертки}
% ------------------------------------------------------------

\begin{examplebox}{Минимальный пример двумерной свертки}

Рассмотрим одноканальный вход
\[
    X=
    \begin{bmatrix}
        1&2&0\\
        0&1&3\\
        2&1&0
    \end{bmatrix}
\]
и ядро
\[
    K=
    \begin{bmatrix}
        1&0\\
        0&-1
    \end{bmatrix}.
\]

Используем взаимную корреляцию без дополнения, с шагом \(1\).
Размер результата:
\[
    H_{\mathrm{out}}=W_{\mathrm{out}}
    =
    \frac{3-2}{1}+1=2.
\]

Первое значение:
\[
    Y_{0,0}
    =
    1\cdot1+2\cdot0+0\cdot0+1\cdot(-1)
    =0.
\]

Второе значение первой строки:
\[
    Y_{0,1}
    =
    2\cdot1+0\cdot0+1\cdot0+3\cdot(-1)
    =-1.
\]

Остальные элементы:
\[
    Y_{1,0}
    =
    0\cdot1+1\cdot0+2\cdot0+1\cdot(-1)
    =-1,
\]
\[
    Y_{1,1}
    =
    1\cdot1+3\cdot0+1\cdot0+0\cdot(-1)
    =1.
\]

Следовательно,
\[
    Y=
    \begin{bmatrix}
        0&-1\\
        -1&1
    \end{bmatrix}.
\]

Если после свертки применяется ReLU, то
\[
    \ReLU(Y)
    =
    \begin{bmatrix}
        0&0\\
        0&1
    \end{bmatrix}.
\]
\end{examplebox}

% ------------------------------------------------------------
\subsection{Число параметров и вычислительная сложность}
% ------------------------------------------------------------

Число параметров стандартного сверточного слоя равно
\[
    N_{\mathrm{params}}
    =
    k_hk_wC_{\mathrm{in}}C_{\mathrm{out}}
    +
    C_{\mathrm{out}},
\]
где последнее слагаемое соответствует смещениям.

Важно, что число параметров не зависит от \(H\) и \(W\). Один фильтр
используется во всех пространственных позициях.

Приближенное число операций умножения-сложения:
\[
    \operatorname{MAC}
    \approx
    H_{\mathrm{out}}W_{\mathrm{out}}
    C_{\mathrm{out}}C_{\mathrm{in}}k_hk_w.
\]

Например, слой с \(C_{\mathrm{in}}=3\), \(C_{\mathrm{out}}=16\) и
ядром \(3\times3\) содержит
\[
    3\cdot3\cdot3\cdot16+16=448
\]
обучаемых параметров.

% ------------------------------------------------------------
\subsection{Эквивариантность относительно сдвига}
% ------------------------------------------------------------

\begin{definition}
Пусть \(T_{\bm{a}}\) --- оператор сдвига:
\[
    (T_{\bm{a}}x)(\bm{n})
    =
    x(\bm{n}-\bm{a}).
\]
Отображение \(F\) называется \emph{эквивариантным} относительно
сдвига, если
\[
    F(T_{\bm{a}}x)
    =
    T_{\bm{a}}F(x).
\]
Оно называется \emph{инвариантным}, если
\[
    F(T_{\bm{a}}x)=F(x).
\]
\end{definition}

Эквивариантность означает, что преобразование входа вызывает
предсказуемое преобразование выхода. Инвариантность означает, что
выход вообще не изменяется.

\begin{theorem}[Эквивариантность свертки относительно сдвига]
Пусть многоканальный сигнал определен на \(\mathbb{Z}^d\), а ядро
имеет конечный носитель. Определим сверточный оператор
\[
    (\mathcal{C}_Kx)_o(\bm{n})
    =
    \sum_c\sum_{\bm{m}}
    K_{o,c}(\bm{m})x_c(\bm{n}+\bm{m}).
\]
Тогда для любого целочисленного сдвига \(\bm{a}\)
\[
    \mathcal{C}_K(T_{\bm{a}}x)
    =
    T_{\bm{a}}(\mathcal{C}_Kx).
\]
\end{theorem}

\begin{proof}
Рассмотрим выходной канал \(o\) и позицию \(\bm{n}\). По определению
свертки
\[
\begin{split}
    \bigl(\mathcal{C}_K(T_{\bm{a}}x)\bigr)_o(\bm{n})
    &=
    \sum_c\sum_{\bm{m}}
    K_{o,c}(\bm{m})
    (T_{\bm{a}}x)_c(\bm{n}+\bm{m}).
\end{split}
\]
По определению сдвига
\[
    (T_{\bm{a}}x)_c(\bm{n}+\bm{m})
    =
    x_c(\bm{n}+\bm{m}-\bm{a}).
\]
Следовательно,
\[
\begin{split}
    \bigl(\mathcal{C}_K(T_{\bm{a}}x)\bigr)_o(\bm{n})
    &=
    \sum_c\sum_{\bm{m}}
    K_{o,c}(\bm{m})
    x_c((\bm{n}-\bm{a})+\bm{m})\\
    &=
    (\mathcal{C}_Kx)_o(\bm{n}-\bm{a})\\
    &=
    \bigl(T_{\bm{a}}(\mathcal{C}_Kx)\bigr)_o(\bm{n}).
\end{split}
\]
Равенство выполнено для всех каналов и позиций, поэтому
\[
    \mathcal{C}_K(T_{\bm{a}}x)
    =
    T_{\bm{a}}(\mathcal{C}_Kx).
\]
\end{proof}

Поточечная нелинейность также коммутирует со сдвигом:
\[
    \varphi(T_{\bm{a}}x)
    =
    T_{\bm{a}}\varphi(x).
\]
Поэтому последовательность сверточных слоев с поточечными
нелинейностями сохраняет эквивариантность.

\begin{importantbox}{Ограничения эквивариантности}
Точная эквивариантность может нарушаться:

\begin{itemize}
    \item на границах из-за дополнения нулями;
    \item при свертке с шагом больше \(1\);
    \item при pooling с шагом больше \(1\);
    \item при округлении координат и изменении размера изображения.
\end{itemize}

Слой с шагом \(s\) в общем случае эквивариантен только относительно
сдвигов, кратных \(s\).
\end{importantbox}

Глобальное усреднение
\[
    \operatorname{GAP}(X)_c
    =
    \frac{1}{HW}
    \sum_{i=1}^{H}\sum_{j=1}^{W}X_{c,i,j}
\]
превращает пространственно эквивариантное представление в
приближенно инвариантное: перестановка пространственных позиций не
меняет их среднее значение.

% ------------------------------------------------------------
\subsection{Рецептивное поле}
% ------------------------------------------------------------

\begin{definition}
\emph{Рецептивное поле} нейрона --- область исходного входа, которая
может повлиять на значение этого нейрона.
\end{definition}

Пусть для слоя \(l\):

\begin{itemize}
    \item \(r_l\) --- размер рецептивного поля;
    \item \(j_l\) --- расстояние во входных пикселях между соседними
    элементами карты признаков;
    \item \(s_l\) --- шаг слоя;
    \item \(k_l^{\mathrm{eff}}\) --- эффективный размер ядра.
\end{itemize}

Начальные значения:
\[
    r_0=1,
    \qquad
    j_0=1.
\]

Рекуррентные формулы:
\[
    j_l=j_{l-1}s_l,
\]
\[
    r_l
    =
    r_{l-1}
    +
    \bigl(k_l^{\mathrm{eff}}-1\bigr)j_{l-1}.
\]

\begin{example}
Рассмотрим последовательность слоев:
\[
    \operatorname{Conv}_{3,s=1}
    \to
    \operatorname{Conv}_{3,s=1}
    \to
    \operatorname{Pool}_{2,s=2}
    \to
    \operatorname{Conv}_{3,s=1}.
\]

Получаем:
\[
\begin{array}{c|c|c}
    \text{слой} & r_l & j_l\\
    \hline
    \text{вход} & 1 & 1\\
    \text{Conv }3\times3 & 3 & 1\\
    \text{Conv }3\times3 & 5 & 1\\
    \text{Pool }2\times2 & 6 & 2\\
    \text{Conv }3\times3 & 10 & 2
\end{array}
\]

Таким образом, элемент последней карты признаков зависит от области
\(10\times10\) исходного изображения.
\end{example}

% ------------------------------------------------------------
\subsection{Функции активации}
% ------------------------------------------------------------

\paragraph{ReLU.}
\[
    \ReLU(x)=\max(0,x).
\]

Преимущества: простая производная, отсутствие насыщения при \(x>0\),
разреженные активации. Недостаток --- возможность появления
неактивных нейронов, для которых вход всегда отрицателен.

\paragraph{Leaky ReLU.}
\[
    \operatorname{LeakyReLU}(x)
    =
    \begin{cases}
        x, & x\geq 0,\\
        \alpha x, & x<0,
    \end{cases}
    \qquad 0<\alpha<1.
\]

Она сохраняет ненулевой градиент в отрицательной полуплоскости и часто
используется в дискриминаторах GAN.

\paragraph{Сигмоида.}
\[
    \sigma(x)=\frac{1}{1+e^{-x}}.
\]
Применяется для получения вероятности в бинарной классификации.
На больших по модулю значениях насыщается, вследствие чего градиент
становится мал.

\paragraph{Гиперболический тангенс.}
\[
    \tanh(x)
    =
    \frac{e^x-e^{-x}}{e^x+e^{-x}}.
\]
Его диапазон равен \((-1,1)\). Поэтому \(\tanh\) часто используется на
выходе генератора, если изображения нормализованы в диапазон
\([-1,1]\).

% ------------------------------------------------------------
\subsection{Pooling и уменьшение разрешения}
% ------------------------------------------------------------

\paragraph{Max pooling.}
Для окна \(\Omega_{i,j}\)
\[
    Y_{c,i,j}
    =
    \max_{(u,v)\in\Omega_{i,j}}X_{c,u,v}.
\]

\paragraph{Average pooling.}
\[
    Y_{c,i,j}
    =
    \frac{1}{|\Omega_{i,j}|}
    \sum_{(u,v)\in\Omega_{i,j}}X_{c,u,v}.
\]

Например, для окна
\[
    \begin{bmatrix}
        1&5\\
        2&3
    \end{bmatrix}
\]
max pooling возвращает \(5\), а average pooling возвращает
\[
    \frac{1+5+2+3}{4}=\frac{11}{4}.
\]

Pooling уменьшает пространственное разрешение и вычислительную
сложность следующих слоев. В современных архитектурах вместо pooling
часто используется свертка с шагом \(2\), параметры которой обучаются.

% ------------------------------------------------------------
\subsection{Нормализация}
% ------------------------------------------------------------

Для mini-batch размера \(B\) и карты \(H\times W\) Batch Normalization
вычисляет статистики отдельно для каждого канала:
\[
    \mu_c
    =
    \frac{1}{BHW}
    \sum_{b,i,j}x_{b,c,i,j},
\]
\[
    \sigma_c^2
    =
    \frac{1}{BHW}
    \sum_{b,i,j}(x_{b,c,i,j}-\mu_c)^2.
\]

Нормированная активация:
\[
    \widehat{x}_{b,c,i,j}
    =
    \frac{x_{b,c,i,j}-\mu_c}
    {\sqrt{\sigma_c^2+\varepsilon}},
\]
а выход:
\[
    y_{b,c,i,j}
    =
    \gamma_c\widehat{x}_{b,c,i,j}+\beta_c.
\]

Параметры \(\gamma_c\) и \(\beta_c\) обучаются. Во время применения
модели используются накопленные оценки среднего и дисперсии.

Для малых batch-размеров могут использоваться Layer Normalization,
Group Normalization или Instance Normalization, статистики которых
слабее зависят или вообще не зависят от других примеров mini-batch.

% ------------------------------------------------------------
\subsection{Специальные виды сверток}
% ------------------------------------------------------------

\subsubsection{Свертка \(1\times1\)}

Свертка \(1\times1\) не объединяет соседние пространственные позиции,
но смешивает каналы:
\[
    Y_{o,i,j}
    =
    b_o+
    \sum_{c=1}^{C_{\mathrm{in}}}
    K_{o,c}X_{c,i,j}.
\]

Она применяется для изменения числа каналов, построения bottleneck-
блоков и снижения вычислительной сложности.

\subsubsection{Групповая свертка}

Входные и выходные каналы разделяются на \(g\) групп. Каждый выходной
канал зависит только от входных каналов своей группы.

Число параметров:
\[
    N_{\mathrm{params}}
    =
    \frac{
        k_hk_wC_{\mathrm{in}}C_{\mathrm{out}}
    }{g}
    +
    C_{\mathrm{out}}.
\]

\subsubsection{Depthwise separable convolution}

Операция состоит из двух частей:

\begin{enumerate}
    \item depthwise-свертка: отдельное пространственное ядро для
    каждого входного канала;
    \item pointwise-свертка \(1\times1\), смешивающая каналы.
\end{enumerate}

Число весов:
\[
    k_hk_wC_{\mathrm{in}}
    +
    C_{\mathrm{in}}C_{\mathrm{out}},
\]
тогда как стандартная свертка требует
\[
    k_hk_wC_{\mathrm{in}}C_{\mathrm{out}}
\]
весов.

\subsubsection{Дилатированная свертка}

Дилатация вставляет интервалы между элементами ядра. Она увеличивает
рецептивное поле без пропорционального увеличения числа параметров.

Для ядра размера \(k\) и дилатации \(d\):
\[
    k^{\mathrm{eff}}=d(k-1)+1.
\]

Например, ядро \(3\times3\) с \(d=2\) имеет эффективный размер
\(5\times5\), но содержит только \(9\) обучаемых коэффициентов на пару
входного и выходного каналов.

\subsubsection{Транспонированная свертка}

Обычную свертку можно представить как умножение на разреженную
матрицу:
\[
    \bm{y}=A\bm{x}.
\]
Транспонированная свертка реализует преобразование
\[
    \bm{z}=A^\top\bm{y}.
\]

Она не является обратной операцией к свертке в общем случае. Название
указывает только на транспонирование матрицы линейного оператора.

Для одного пространственного измерения размер выхода равен
\[
    H_{\mathrm{out}}
    =
    (H_{\mathrm{in}}-1)s
    -2p
    +d(k-1)
    +o
    +1,
\]
где \(o\) --- параметр \texttt{output padding}.

Транспонированные свертки используются в декодерах и генераторах.
Неравномерное перекрытие ядер может создавать шахматные артефакты.
Альтернативой служит явное увеличение разрешения, например
билинейная интерполяция, с последующей обычной сверткой.

% ------------------------------------------------------------
\subsection{Остаточные связи}
% ------------------------------------------------------------

Остаточный блок имеет вид
\[
    \bm{y}
    =
    \varphi\bigl(\mathcal{F}(\bm{x};\theta)+\bm{x}\bigr).
\]

Если размерности различаются, используется проекция:
\[
    \bm{y}
    =
    \varphi\bigl(
        \mathcal{F}(\bm{x};\theta)
        +
        W_s\bm{x}
    \bigr).
\]

Для отображения
\[
    \bm{y}=\mathcal{F}(\bm{x})+\bm{x}
\]
якобиан равен
\[
    \frac{\partial\bm{y}}{\partial\bm{x}}
    =
    \frac{\partial\mathcal{F}}{\partial\bm{x}}+I.
\]
Слагаемое \(I\) создает прямой путь распространения градиента и
облегчает обучение глубоких сетей.

% ------------------------------------------------------------
\subsection{Выходные слои и функции потерь}
% ------------------------------------------------------------

В задаче классификации на \(C\) классов сеть выдает логиты
\(\bm{a}\in\R^C\). Вероятности определяются softmax:
\[
    p_c
    =
    \frac{e^{a_c}}{\sum_{j=1}^{C}e^{a_j}}.
\]

Кросс-энтропия:
\[
    \ell(\bm{a},\bm{y})
    =
    -\sum_{c=1}^{C}y_c\log p_c,
\]
где \(\sum_c y_c=1\).

\begin{proposition}[Градиент softmax-кросс-энтропии]
Пусть
\[
    p_c=\frac{e^{a_c}}{\sum_j e^{a_j}},
    \qquad
    \ell=-\sum_c y_c\log p_c,
    \qquad
    \sum_c y_c=1.
\]
Тогда
\[
    \frac{\partial\ell}{\partial a_j}
    =
    p_j-y_j.
\]
\end{proposition}

\begin{proof}
Запишем логарифм softmax:
\[
    \log p_c
    =
    a_c-\log\left(\sum_k e^{a_k}\right).
\]
Тогда
\[
\begin{split}
    \ell
    &=
    -\sum_c y_c a_c
    +
    \sum_c y_c
    \log\left(\sum_k e^{a_k}\right)\\
    &=
    -\sum_c y_c a_c
    +
    \log\left(\sum_k e^{a_k}\right),
\end{split}
\]
поскольку \(\sum_c y_c=1\).

Дифференцируя по \(a_j\), получаем
\[
\begin{split}
    \frac{\partial\ell}{\partial a_j}
    &=
    -y_j
    +
    \frac{e^{a_j}}{\sum_k e^{a_k}}\\
    &=
    p_j-y_j.
\end{split}
\]
\end{proof}

В бинарной классификации используется бинарная кросс-энтропия:
\[
    \ell(s,y)
    =
    -y\log\sigma(s)
    -(1-y)\log(1-\sigma(s)).
\]

Для численной устойчивости функции потерь обычно вычисляются
непосредственно по логитам, без явного вычисления логарифмов
вероятностей, близких к нулю.

% ------------------------------------------------------------
\subsection{Обратное распространение через свертку}
% ------------------------------------------------------------

Рассмотрим простой случай: единичный шаг, отсутствие дополнения и
дилатации:
\[
    Y_{o,i,j}
    =
    b_o+
    \sum_{c,u,v}
    K_{o,c,u,v}X_{c,i+u,j+v}.
\]

Обозначим
\[
    \Delta_{o,i,j}
    =
    \frac{\partial L}{\partial Y_{o,i,j}}.
\]

Тогда градиент по ядру:
\[
    \frac{\partial L}{\partial K_{o,c,u,v}}
    =
    \sum_{i,j}
    \Delta_{o,i,j}X_{c,i+u,j+v}.
\]

Градиент по смещению:
\[
    \frac{\partial L}{\partial b_o}
    =
    \sum_{i,j}\Delta_{o,i,j}.
\]

Градиент по входу:
\[
    \frac{\partial L}{\partial X_{c,p,q}}
    =
    \sum_{o,u,v}
    K_{o,c,u,v}
    \Delta_{o,p-u,q-v},
\]
где значения \(\Delta\) за пределами допустимой области полагаются
равными нулю.

Таким образом, вычисление градиентов также выражается через операции,
близкие к свертке и транспонированной свертке.

% ------------------------------------------------------------
\subsection{Алгоритм обучения CNN}
% ------------------------------------------------------------

\begin{algobox}{Алгоритм обучения сверточной сети}
\textbf{Вход:} обучающая выборка, сеть \(f_\theta\), функция потерь,
размер mini-batch \(B\), оптимизатор.

\begin{enumerate}
    \item Инициализировать параметры \(\theta\).

    \item Разделить данные на обучающую, валидационную и тестовую
    части.

    \item Для каждой эпохи:
    \begin{enumerate}
        \item перемешать обучающие объекты;
        \item разделить выборку на mini-batch;
        \item для каждого mini-batch:
        \begin{enumerate}
            \item выполнить преобразование и аугментацию данных;
            \item вычислить прямой проход
            \[
                \widehat{\bm{y}}_i=f_\theta(\bm{x}_i);
            \]
            \item вычислить функцию потерь
            \[
                L(\theta)
                =
                \frac{1}{B}\sum_{i=1}^{B}
                \ell(\widehat{\bm{y}}_i,\bm{y}_i)
                +
                \frac{\lambda}{2}\|\theta\|_2^2;
            \]
            \item выполнить обратное распространение и найти
            \(g=\nabla_\theta L\);
            \item обновить параметры оптимизатором;
            \item обнулить накопленные градиенты.
        \end{enumerate}
        \item оценить качество на валидационной выборке;
        \item при необходимости изменить скорость обучения или
        применить раннюю остановку.
    \end{enumerate}

    \item Один раз оценить итоговую модель на тестовой выборке.
\end{enumerate}
\end{algobox}

\subsubsection{Оптимизатор Adam}

Пусть \(g_t=\nabla_\theta L_t\) --- стохастический градиент на шаге
\(t\). Adam поддерживает экспоненциальные средние первого и второго
моментов:
\[
    m_t
    =
    \beta_1m_{t-1}
    +(1-\beta_1)g_t,
\]
\[
    v_t
    =
    \beta_2v_{t-1}
    +(1-\beta_2)g_t^{\odot2},
\]
где квадрат вычисляется покомпонентно.

Коррекция смещения:
\[
    \widehat{m}_t
    =
    \frac{m_t}{1-\beta_1^t},
    \qquad
    \widehat{v}_t
    =
    \frac{v_t}{1-\beta_2^t}.
\]

Обновление:
\[
    \theta_t
    =
    \theta_{t-1}
    -
    \eta
    \frac{\widehat{m}_t}
    {\sqrt{\widehat{v}_t}+\varepsilon}.
\]

В GAN для генератора и дискриминатора используются отдельные
экземпляры оптимизатора с независимыми состояниями \(m_t,v_t\).

% ------------------------------------------------------------
\subsection{Регуляризация CNN}
% ------------------------------------------------------------

Основные методы регуляризации:

\begin{enumerate}
    \item \textbf{Аугментация данных:} отражения, повороты, случайные
    обрезки, изменение яркости и контраста.

    \item \textbf{Weight decay:}
    \[
        L_{\mathrm{reg}}
        =
        L+\frac{\lambda}{2}\|\theta\|_2^2.
    \]

    \item \textbf{Dropout:}
    \[
        \widetilde{\bm{h}}
        =
        \frac{\bm{m}\odot\bm{h}}{1-p},
        \qquad
        m_i\sim\operatorname{Bernoulli}(1-p).
    \]
    При применении модели случайная маска не используется.

    \item \textbf{Ранняя остановка} по валидационной метрике.

    \item \textbf{Сглаживание меток:}
    вместо строго one-hot-вектора используется распределение
    \[
        y'_c
        =
        (1-\varepsilon)y_c+\frac{\varepsilon}{C}.
    \]

    \item \textbf{Уменьшение числа параметров:} global average
    pooling, групповые и depthwise-свертки.
\end{enumerate}

Для ReLU-сетей веса часто инициализируют случайными значениями с
дисперсией порядка
\[
    \operatorname{Var}(W_{ij})
    =
    \frac{2}{\operatorname{fan}_{\mathrm{in}}}.
\]

% ------------------------------------------------------------
\subsection{Минимальный пример архитектуры CNN}
% ------------------------------------------------------------

Пусть требуется классифицировать изображения
\(28\times28\times1\) на \(10\) классов.

\[
\begin{array}{c|c|c}
    \text{слой}
    &
    \text{размер выхода}
    &
    \text{число параметров}\\
    \hline
    \text{вход}
    &
    1\times28\times28
    &
    0\\
    \text{Conv }3\times3,\;1\to8,\;\text{same}
    &
    8\times28\times28
    &
    3\cdot3\cdot1\cdot8+8=80\\
    \text{ReLU}
    &
    8\times28\times28
    &
    0\\
    \text{MaxPool }2\times2
    &
    8\times14\times14
    &
    0\\
    \text{Conv }3\times3,\;8\to16,\;\text{same}
    &
    16\times14\times14
    &
    3\cdot3\cdot8\cdot16+16=1168\\
    \text{ReLU}
    &
    16\times14\times14
    &
    0\\
    \text{Global Average Pooling}
    &
    16
    &
    0\\
    \text{Linear }16\to10
    &
    10
    &
    16\cdot10+10=170
\end{array}
\]

Общее число параметров:
\[
    80+1168+170=1418.
\]

После получения логитов применяется softmax-кросс-энтропия. Softmax
необязательно включать в саму сеть, если функция потерь принимает
логиты.

% ------------------------------------------------------------
\subsection{Основные архитектурные семейства}
% ------------------------------------------------------------

\begin{center}
\begin{tabularx}{\textwidth}{
    >{\bfseries}p{2.5cm}
    X
    X
}
\toprule
Архитектура & Основная идея & Характерные особенности\\
\midrule
LeNet
&
Раннее применение обучаемых сверток для распознавания символов.
&
Свертки, subsampling, полносвязный классификатор.\\

AlexNet
&
Глубокая CNN для крупномасштабной классификации.
&
ReLU, GPU-обучение, аугментация, dropout, max pooling.\\

VGG
&
Увеличение глубины при использовании малых ядер.
&
Последовательности сверток \(3\times3\).\\

ResNet
&
Обучение остаточного отображения.
&
Skip connections и глубокие остаточные блоки.\\

Fully convolutional network
&
Замена полносвязной головы свертками.
&
Пространственный выход переменного размера.\\

U-Net
&
Encoder--decoder с передачей детальных признаков.
&
Skip connections между уровнями одинакового разрешения.\\
\bottomrule
\end{tabularx}
\end{center}

% ------------------------------------------------------------
\subsection{Основные задачи CNN}
% ------------------------------------------------------------

\begin{enumerate}
    \item \textbf{Классификация:}
    один вектор вероятностей для всего изображения.

    \item \textbf{Сегментация:}
    предсказание класса для каждого пикселя.

    \item \textbf{Детектирование:}
    предсказание классов и ограничивающих прямоугольников.

    \item \textbf{Оценка ключевых точек:}
    построение карт вероятностей положения ориентиров.

    \item \textbf{Восстановление изображений:}
    удаление шума, super-resolution, устранение размытия.

    \item \textbf{Генерация:}
    использование сверточных декодеров в GAN, автоэнкодерах и других
    генеративных моделях.

    \item \textbf{Одномерные и трехмерные данные:}
    аудиосигналы, временные ряды, томографические изображения и видео.
\end{enumerate}

% ------------------------------------------------------------
\subsection{Типичные ошибки при проектировании CNN}
% ------------------------------------------------------------

\begin{enumerate}
    \item Неверный расчет пространственных размеров после свертки и
    pooling.

    \item Слишком раннее уменьшение разрешения, приводящее к потере
    мелких объектов.

    \item Использование большой полносвязной головы после
    \texttt{flatten}, создающей чрезмерное число параметров.

    \item Несогласованная нормализация обучающих и тестовых данных.

    \item Утечка тестовых данных при настройке гиперпараметров.

    \item Оценка только accuracy при сильном дисбалансе классов.

    \item Использование Batch Normalization с крайне малым batch без
    учета нестабильности статистик.

    \item Предположение о точной инвариантности CNN к сдвигам:
    padding, stride и pooling могут ее нарушать.
\end{enumerate}

% ============================================================
\section{Генеративные состязательные сети}
% ============================================================

\subsection{Генеративное моделирование}

Пусть реальные данные имеют неизвестное распределение
\(P_{\mathrm{data}}\). Генеративная модель должна построить
распределение \(P_g\), близкое к \(P_{\mathrm{data}}\), и позволять
получать новые выборки
\[
    \widetilde{\bm{x}}\sim P_g.
\]

В GAN распределение задается неявно. Сначала выбирается простое
латентное распределение
\[
    \bm{z}\sim P_z,
\]
например
\[
    \bm{z}\sim\mathcal{N}(\bm{0},I)
\]
или равномерное распределение. Затем генератор вычисляет
\[
    \widetilde{\bm{x}}=G_\theta(\bm{z}).
\]

Индуцированное генератором распределение называется pushforward:
\[
    P_g=G_{\theta\#}P_z.
\]
Для измеримого множества \(A\)
\[
    P_g(A)
    =
    P_z\bigl(G_\theta^{-1}(A)\bigr).
\]

Вычислять плотность \(p_g(\bm{x})\) при этом обычно не требуется.

% ------------------------------------------------------------
\subsection{Компоненты GAN}
% ------------------------------------------------------------

GAN состоит из двух моделей:

\begin{enumerate}
    \item \textbf{Генератор}
    \[
        G_\theta:\mathcal{Z}\to\mathcal{X},
    \]
    преобразующий латентный шум в синтетический объект.

    \item \textbf{Дискриминатор}
    \[
        D_\phi:\mathcal{X}\to[0,1],
    \]
    оценивающий вероятность того, что объект является реальным.
\end{enumerate}

\begin{center}
\begin{tikzpicture}[node distance=10mm]
    \node[block] (z) {Шум \(\bm{z}\)};
    \node[block,right=of z] (g) {Генератор\\\(G_\theta\)};
    \node[block,right=of g] (fake) {Синтетический\\объект};
    \node[block,right=15mm of fake] (d) {Дискриминатор\\\(D_\phi\)};
    \node[block,above=10mm of fake] (real) {Реальный\\объект};
    \node[block,right=of d] (score) {Оценка\\реальности};

    \draw[arrow] (z) -- (g);
    \draw[arrow] (g) -- (fake);
    \draw[arrow] (fake) -- (d);
    \draw[arrow] (real) -| (d);
    \draw[arrow] (d) -- (score);
\end{tikzpicture}
\end{center}

Дискриминатор обучается различать реальные и синтетические данные.
Генератор обучается создавать объекты, которые дискриминатор
классифицирует как реальные.

% ------------------------------------------------------------
\subsection{Минимаксная функция GAN}
% ------------------------------------------------------------

Классическая функция ценности:
\[
\begin{split}
    V(D,G)
    &=
    \E_{\bm{x}\sim P_{\mathrm{data}}}
    [\log D(\bm{x})]\\
    &\quad+
    \E_{\bm{z}\sim P_z}
    [\log(1-D(G(\bm{z})))].
\end{split}
\]

Игра задается задачей
\[
    \min_G\max_D V(D,G).
\]

Дискриминатор максимизирует вероятность правильной классификации.
Генератор минимизирует ту часть критерия, которая относится к
синтетическим данным.

Для фиксированного генератора максимизация по \(D\) эквивалентна
обучению бинарного классификатора:

\begin{itemize}
    \item реальные данные имеют метку \(1\);
    \item синтетические данные имеют метку \(0\).
\end{itemize}

% ------------------------------------------------------------
\subsection{Неотрицательность KL-дивергенции}
% ------------------------------------------------------------

Для дальнейшего анализа потребуется следующее утверждение.

\begin{lemma}[Неравенство Гиббса]
Пусть \(P\) и \(Q\) --- вероятностные распределения с плотностями
\(p\) и \(q\) относительно общей меры. Тогда
\[
    \KL(P\|Q)
    =
    \int p(x)\log\frac{p(x)}{q(x)}\,dx
    \geq 0.
\]
Если дивергенция конечна, равенство достигается тогда и только тогда,
когда \(p=q\) почти всюду.
\end{lemma}

\begin{proof}
Используем неравенство
\[
    \log t\leq t-1,
    \qquad t>0,
\]
причем равенство достигается только при \(t=1\).

На множестве, где \(p(x)>0\) и \(q(x)>0\), подставим
\[
    t=\frac{q(x)}{p(x)}.
\]
Тогда
\[
    \log\frac{q(x)}{p(x)}
    \leq
    \frac{q(x)}{p(x)}-1.
\]
Умножая на \(p(x)\), получаем
\[
    p(x)\log\frac{q(x)}{p(x)}
    \leq
    q(x)-p(x).
\]
Интегрирование дает
\[
    \int p(x)\log\frac{q(x)}{p(x)}\,dx
    \leq
    \int q(x)\,dx-\int p(x)\,dx=0.
\]
Следовательно,
\[
    -\int p(x)\log\frac{q(x)}{p(x)}\,dx
    =
    \int p(x)\log\frac{p(x)}{q(x)}\,dx
    \geq 0.
\]

Если существует множество положительной \(P\)-меры, на котором
\(q=0\), то \(\KL(P\|Q)=+\infty\), и неравенство также выполнено.

Равенство в использованном неравенстве возможно только при
\[
    \frac{q(x)}{p(x)}=1
\]
почти всюду относительно \(P\). С учетом нормировки распределений это
означает \(p=q\) почти всюду. Обратное утверждение очевидно:
если \(p=q\), то подынтегральное выражение равно нулю.
\end{proof}

% ------------------------------------------------------------
\subsection{Оптимальный дискриминатор}
% ------------------------------------------------------------

\begin{theorem}[Оптимальный дискриминатор для фиксированного генератора]
Пусть \(P_{\mathrm{data}}\) и \(P_g\) имеют плотности
\(p_{\mathrm{data}}\) и \(p_g\) относительно общей меры. Для
фиксированного генератора максимум функционала
\[
    V(D,G)
    =
    \int
    \left[
        p_{\mathrm{data}}(x)\log D(x)
        +
        p_g(x)\log(1-D(x))
    \right]dx
\]
по измеримым функциям \(D:\mathcal{X}\to[0,1]\) достигается при
\[
    D^*(x)
    =
    \frac{
        p_{\mathrm{data}}(x)
    }{
        p_{\mathrm{data}}(x)+p_g(x)
    }
\]
для почти всех \(x\), где
\(p_{\mathrm{data}}(x)+p_g(x)>0\). В точках, где обе плотности равны
нулю, значение \(D(x)\) произвольно.
\end{theorem}

\begin{proof}
Зафиксируем точку \(x\) и обозначим
\[
    a=p_{\mathrm{data}}(x),
    \qquad
    b=p_g(x).
\]
Требуется максимизировать по \(d\in[0,1]\) функцию
\[
    f(d)=a\log d+b\log(1-d).
\]

Пусть сначала \(a>0\) и \(b>0\). Тогда
\[
    f'(d)
    =
    \frac{a}{d}-\frac{b}{1-d}.
\]
Условие стационарности \(f'(d)=0\) дает
\[
    \frac{a}{d}
    =
    \frac{b}{1-d}.
\]
Отсюда
\[
    a(1-d)=bd,
\]
\[
    a=(a+b)d,
\]
\[
    d=\frac{a}{a+b}.
\]

Вторая производная:
\[
    f''(d)
    =
    -\frac{a}{d^2}
    -
    \frac{b}{(1-d)^2}<0.
\]
Следовательно, функция строго вогнута, а найденная стационарная точка
является единственным максимумом.

Если \(a>0\), \(b=0\), максимизируется \(a\log d\), поэтому максимум
достигается при \(d=1\). Формула также дает
\[
    d=\frac{a}{a}=1.
\]

Если \(a=0\), \(b>0\), максимизируется \(b\log(1-d)\), поэтому
\(d=0\). Формула дает
\[
    d=\frac{0}{b}=0.
\]

Если \(a=b=0\), значение подынтегрального выражения равно нулю для
любого \(d\).

Таким образом, поточечный максимум достигается при
\[
    D^*(x)
    =
    \frac{
        p_{\mathrm{data}}(x)
    }{
        p_{\mathrm{data}}(x)+p_g(x)
    }.
\]
Поскольку интеграл представляет собой сумму независимых поточечных
слагаемых, выбор поточечного максимизатора максимизирует весь
функционал.
\end{proof}

Оптимальный дискриминатор оценивает отношение плотностей. В частности,
\[
    \frac{D^*(x)}{1-D^*(x)}
    =
    \frac{p_{\mathrm{data}}(x)}{p_g(x)}
\]
в точках, где обе плотности положительны.

% ------------------------------------------------------------
\subsection{Глобальный оптимум классической GAN}
% ------------------------------------------------------------

Определим смесь
\[
    m(x)
    =
    \frac{
        p_{\mathrm{data}}(x)+p_g(x)
    }{2}.
\]

Дивергенция Дженсена--Шеннона:
\[
\begin{split}
    \JS(P_{\mathrm{data}}\|P_g)
    &=
    \frac{1}{2}
    \KL(P_{\mathrm{data}}\|M)\\
    &\quad+
    \frac{1}{2}
    \KL(P_g\|M),
\end{split}
\]
где \(M\) имеет плотность \(m\).

\begin{theorem}[Глобальный оптимум классической GAN]
Пусть дискриминатор может представлять любую измеримую функцию
\(D:\mathcal{X}\to[0,1]\). Для фиксированного генератора определим
\[
    C(G)=\max_D V(D,G).
\]
Тогда
\[
    C(G)
    =
    -\log4
    +
    2\JS(P_{\mathrm{data}}\|P_g).
\]
Следовательно,
\[
    C(G)\geq-\log4,
\]
и равенство достигается тогда и только тогда, когда
\[
    P_g=P_{\mathrm{data}}.
\]
\end{theorem}

\begin{proof}
Подставим оптимальный дискриминатор:
\[
\begin{split}
    C(G)
    &=
    \int
    p_{\mathrm{data}}(x)
    \log
    \frac{
        p_{\mathrm{data}}(x)
    }{
        p_{\mathrm{data}}(x)+p_g(x)
    }
    dx\\
    &\quad+
    \int
    p_g(x)
    \log
    \frac{
        p_g(x)
    }{
        p_{\mathrm{data}}(x)+p_g(x)
    }
    dx.
\end{split}
\]

Поскольку
\[
    p_{\mathrm{data}}(x)+p_g(x)=2m(x),
\]
получаем
\[
\begin{split}
    C(G)
    &=
    \int
    p_{\mathrm{data}}(x)
    \log
    \frac{
        p_{\mathrm{data}}(x)
    }{2m(x)}
    dx\\
    &\quad+
    \int
    p_g(x)
    \log
    \frac{
        p_g(x)
    }{2m(x)}
    dx.
\end{split}
\]

Разделим логарифмы:
\[
    \log\frac{p(x)}{2m(x)}
    =
    \log\frac{p(x)}{m(x)}-\log2.
\]
Поэтому
\[
\begin{split}
    C(G)
    &=
    \int
    p_{\mathrm{data}}(x)
    \log
    \frac{
        p_{\mathrm{data}}(x)
    }{m(x)}
    dx\\
    &\quad+
    \int
    p_g(x)
    \log
    \frac{
        p_g(x)
    }{m(x)}
    dx\\
    &\quad-
    \log2\int p_{\mathrm{data}}(x)\,dx
    -
    \log2\int p_g(x)\,dx.
\end{split}
\]

Обе плотности нормированы, поэтому
\[
    \int p_{\mathrm{data}}(x)\,dx
    =
    \int p_g(x)\,dx
    =1.
\]
Следовательно,
\[
\begin{split}
    C(G)
    &=
    \KL(P_{\mathrm{data}}\|M)
    +
    \KL(P_g\|M)
    -
    2\log2\\
    &=
    2\JS(P_{\mathrm{data}}\|P_g)-\log4.
\end{split}
\]

По неравенству Гиббса обе KL-дивергенции неотрицательны. Поэтому
\[
    \JS(P_{\mathrm{data}}\|P_g)\geq0
\]
и
\[
    C(G)\geq-\log4.
\]

Если \(P_g=P_{\mathrm{data}}\), то \(M=P_g=P_{\mathrm{data}}\), обе
KL-дивергенции равны нулю и
\[
    C(G)=-\log4.
\]

Обратно, если
\[
    C(G)=-\log4,
\]
то
\[
    \JS(P_{\mathrm{data}}\|P_g)=0.
\]
Сумма двух неотрицательных KL-дивергенций может быть равна нулю только
тогда, когда каждая из них равна нулю:
\[
    P_{\mathrm{data}}=M,
    \qquad
    P_g=M.
\]
Следовательно,
\[
    P_g=P_{\mathrm{data}}.
\]
\end{proof}

\begin{corollary}
Если \(P_g=P_{\mathrm{data}}\), то оптимальный дискриминатор равен
\[
    D^*(x)=\frac{1}{2}
\]
почти всюду на носителе распределения данных.
\end{corollary}

Важно, что единственным является оптимальное \emph{распределение}
генератора, но не обязательно его параметры. Разные значения
\(\theta\) могут задавать одно и то же распределение \(P_g\).

% ------------------------------------------------------------
\subsection{Насыщающаяся и ненасыщающаяся функция потерь}
% ------------------------------------------------------------

Исходная минимаксная функция потерь генератора:
\[
    L_G^{\mathrm{minimax}}
    =
    \E_{\bm{z}}
    \left[
        \log(1-D(G(\bm{z})))
    \right].
\]

В начале обучения дискриминатор часто легко обнаруживает синтетические
объекты, поэтому
\[
    D(G(\bm{z}))\approx0.
\]
В этом случае градиент минимаксного критерия может быть мал.

Пусть \(a\) --- логит дискриминатора для синтетического объекта:
\[
    D=\sigma(a).
\]
Тогда
\[
    \frac{d}{da}\log(1-\sigma(a))
    =
    -\sigma(a)
    =
    -D.
\]
Если \(D\approx0\), производная близка к нулю.

На практике генератор часто обучают с ненасыщающейся функцией:
\[
    L_G^{\mathrm{NS}}
    =
    -\E_{\bm{z}}
    [\log D(G(\bm{z}))].
\]
Для нее
\[
    \frac{d}{da}\bigl(-\log\sigma(a)\bigr)
    =
    -(1-\sigma(a))
    =
    D-1.
\]
Если \(D\approx0\), производная близка к \(-1\), поэтому генератор
получает более сильный обучающий сигнал.

Ненасыщающаяся функция имеет тот же желательный оптимум
\(P_g=P_{\mathrm{data}}\), но дает другую динамику оптимизации.

% ------------------------------------------------------------
\subsection{Алгоритм обучения классической GAN}
% ------------------------------------------------------------

На практике дискриминатор часто выдает логит \(s_\phi(x)\), а не
вероятность. Тогда
\[
    D_\phi(x)=\sigma(s_\phi(x)).
\]

Численно устойчивые функции потерь:
\[
\begin{split}
    L_D
    &=
    \E_{x\sim P_{\mathrm{data}}}
    [\operatorname{softplus}(-s_\phi(x))]\\
    &\quad+
    \E_{z\sim P_z}
    [\operatorname{softplus}(s_\phi(G_\theta(z)))],
\end{split}
\]
\[
    L_G
    =
    \E_{z\sim P_z}
    [\operatorname{softplus}(-s_\phi(G_\theta(z)))],
\]
где
\[
    \operatorname{softplus}(t)=\log(1+e^t).
\]

\begin{algobox}{Алгоритм обучения ненасыщающейся GAN}
\textbf{Вход:} генератор \(G_\theta\), дискриминатор \(D_\phi\),
распределение шума \(P_z\), число обновлений дискриминатора \(k\).

\begin{enumerate}
    \item Инициализировать \(\theta\) и \(\phi\).

    \item Пока не выполнено условие остановки:
    \begin{enumerate}
        \item Повторить \(k\) раз обновление дискриминатора:
        \begin{enumerate}
            \item выбрать реальные объекты
            \[
                \{x_i\}_{i=1}^{B}
                \sim P_{\mathrm{data}};
            \]
            \item выбрать шум
            \[
                \{z_i\}_{i=1}^{B}
                \sim P_z;
            \]
            \item построить синтетические объекты
            \[
                \widetilde{x}_i=G_\theta(z_i);
            \]
            \item при обновлении дискриминатора не передавать градиент
            в генератор;
            \item вычислить
            \[
            \begin{split}
                \widehat{L}_D
                &=
                -\frac{1}{B}\sum_{i=1}^{B}
                \bigl[
                    \log D_\phi(x_i)\\
                &\hspace{27mm}
                    +\log(1-D_\phi(\widetilde{x}_i))
                \bigr];
            \end{split}
            \]
            \item обновить
            \[
                \phi
                \leftarrow
                \phi-\eta_D\nabla_\phi\widehat{L}_D.
            \]
        \end{enumerate}

        \item Обновить генератор:
        \begin{enumerate}
            \item выбрать новый шум
            \[
                \{z_i\}_{i=1}^{B}\sim P_z;
            \]
            \item вычислить
            \[
                \widehat{L}_G
                =
                -\frac{1}{B}
                \sum_{i=1}^{B}
                \log D_\phi(G_\theta(z_i));
            \]
            \item обновить
            \[
                \theta
                \leftarrow
                \theta-\eta_G\nabla_\theta\widehat{L}_G.
            \]
        \end{enumerate}

        \item Периодически сохранять параметры, фиксированные наборы
        сгенерированных примеров и метрики качества.
    \end{enumerate}
\end{enumerate}
\end{algobox}

Генератор и дискриминатор не минимизируют одну и ту же функцию обычным
градиентным спуском. Один участник игры улучшает способность отличать
распределения, а другой изменяет само генерируемое распределение.

% ------------------------------------------------------------
\subsection{Минимальный дискретный пример}
% ------------------------------------------------------------

\begin{example}
Пусть пространство состоит из двух точек:
\[
    \mathcal{X}=\{a,b\}.
\]
Реальное распределение:
\[
    p_{\mathrm{data}}(a)=\frac{3}{4},
    \qquad
    p_{\mathrm{data}}(b)=\frac{1}{4}.
\]
Распределение генератора:
\[
    p_g(a)=\frac{1}{2},
    \qquad
    p_g(b)=\frac{1}{2}.
\]

Оптимальный дискриминатор:
\[
    D^*(a)
    =
    \frac{
        3/4
    }{
        3/4+1/2
    }
    =
    \frac{3}{5},
\]
\[
    D^*(b)
    =
    \frac{
        1/4
    }{
        1/4+1/2
    }
    =
    \frac{1}{3}.
\]

Точка \(a\) чаще встречается в реальных данных, поэтому
\(D^*(a)>1/2\). Точка \(b\) относительно чаще генерируется моделью,
поэтому \(D^*(b)<1/2\).

Если генератор достигнет
\[
    p_g(a)=\frac{3}{4},
    \qquad
    p_g(b)=\frac{1}{4},
\]
то
\[
    D^*(a)=D^*(b)=\frac{1}{2}.
\]
\end{example}

% ------------------------------------------------------------
\subsection{Минимальный непрерывный пример}
% ------------------------------------------------------------

Пусть реальные данные распределены как
\[
    x\sim\mathcal{N}(2,0.5^2),
\]
а латентная переменная:
\[
    z\sim\mathcal{N}(0,1).
\]

Рассмотрим линейный генератор
\[
    G_{a,b}(z)=az+b.
\]
Тогда
\[
    G_{a,b}(z)\sim\mathcal{N}(b,a^2).
\]

Для точного совпадения распределений необходимо
\[
    b=2,
    \qquad
    a^2=0.25,
\]
то есть
\[
    b=2,
    \qquad
    |a|=0.5.
\]

Это показывает, что разные параметры \(a=0.5\) и \(a=-0.5\) задают
одно и то же оптимальное распределение генератора.

% ------------------------------------------------------------
\subsection{Deep Convolutional GAN}
% ------------------------------------------------------------

Для генерации изображений генератор и дискриминатор обычно строятся на
основе сверточных слоев.

Пример генератора для изображений \(32\times32\times3\):
\[
    z\in\R^{100}
    \longrightarrow
    256\times4\times4
    \longrightarrow
    128\times8\times8
    \longrightarrow
    64\times16\times16
    \longrightarrow
    3\times32\times32.
\]

Первое преобразование может быть полносвязным, а последующие ---
транспонированными свертками или операциями
\[
    \text{upsampling}\to\text{обычная свертка}.
\]

Типичный генератор:

\begin{enumerate}
    \item преобразует \(z\) в небольшую пространственную карту с
    большим числом каналов;
    \item постепенно увеличивает разрешение;
    \item уменьшает число каналов;
    \item использует ReLU во внутренних слоях;
    \item использует \(\tanh\) на выходе для данных в диапазоне
    \([-1,1]\).
\end{enumerate}

Дискриминатор выполняет обратное преобразование:
\[
    3\times32\times32
    \longrightarrow
    64\times16\times16
    \longrightarrow
    128\times8\times8
    \longrightarrow
    256\times4\times4
    \longrightarrow
    \text{логит}.
\]

Типичный дискриминатор:

\begin{enumerate}
    \item использует свертки с шагом \(2\) вместо pooling;
    \item постепенно уменьшает разрешение и увеличивает число каналов;
    \item использует Leaky ReLU;
    \item завершает работу одним скалярным логитом;
    \item не применяет softmax, так как решается бинарная задача.
\end{enumerate}

% ------------------------------------------------------------
\subsection{Основные проблемы обучения GAN}
% ------------------------------------------------------------

\subsubsection{Коллапс мод}

При \emph{mode collapse} разные латентные векторы отображаются в
одинаковые или очень похожие объекты:
\[
    G(z_1)\approx G(z_2)
\]
для большого множества различных \(z_1,z_2\).

Генератор может производить реалистичные объекты, но покрывать только
малую часть распределения данных.

\subsubsection{Слишком сильный дискриминатор}

Если дискриминатор почти идеально разделяет реальные и синтетические
объекты, генератор может получать слабые или плохо направленные
градиенты.

Однако слишком слабый дискриминатор также вреден: его градиенты не
содержат полезной информации о различиях распределений.

\subsubsection{Колебания и отсутствие сходимости}

GAN является динамической игрой. Градиентные методы могут двигаться
по циклическим траекториям вместо приближения к равновесию.

Простейший пример:
\[
    \min_x\max_y xy.
\]
Градиенты:
\[
    \frac{\partial(xy)}{\partial x}=y,
    \qquad
    \frac{\partial(xy)}{\partial y}=x.
\]
Одновременные обновления
\[
    x_{t+1}=x_t-\eta y_t,
    \qquad
    y_{t+1}=y_t+\eta x_t
\]
могут вращаться вокруг точки \((0,0)\), а не сходиться к ней.

\subsubsection{Шахматные артефакты}

Транспонированная свертка может создавать периодические артефакты
из-за неравномерного перекрытия ядер. Возможные решения:

\begin{itemize}
    \item выбирать согласованные размер ядра и шаг;
    \item использовать upsampling с последующей обычной сверткой;
    \item контролировать дополнение и размер выхода.
\end{itemize}

\subsubsection{Переобучение и запоминание}

Генератор может воспроизводить близкие копии обучающих объектов.
Поэтому визуальная реалистичность недостаточна: необходимо проверять
разнообразие и ближайших соседей среди обучающих примеров.

% ------------------------------------------------------------
\subsection{Стабилизация GAN}
% ------------------------------------------------------------

Основные методы:

\begin{enumerate}
    \item ненасыщающаяся функция потерь генератора;

    \item различные скорости обучения для генератора и
    дискриминатора;

    \item несколько обновлений дискриминатора на одно обновление
    генератора;

    \item спектральная нормализация дискриминатора;

    \item штраф на норму градиента;

    \item ограничение величины градиента;

    \item сглаживание реальных меток, например замена \(1\) на
    значение немного меньше \(1\);

    \item аугментация данных на входе дискриминатора;

    \item экспоненциальное среднее параметров генератора:
    \[
        \theta_{\mathrm{EMA}}
        \leftarrow
        \beta\theta_{\mathrm{EMA}}
        +(1-\beta)\theta;
    \]

    \item фиксированный набор латентных векторов для наблюдения
    динамики обучения;

    \item проверка не только качества, но и разнообразия образцов.
\end{enumerate}

% ------------------------------------------------------------
\subsection{Спектральная нормализация}
% ------------------------------------------------------------

\begin{definition}
Спектральная норма матрицы \(W\) определяется как
\[
    \|W\|_2
    =
    \sup_{\bm{x}\neq0}
    \frac{\|W\bm{x}\|_2}{\|\bm{x}\|_2}.
\]
Она равна наибольшему сингулярному числу
\(\sigma_{\max}(W)\).
\end{definition}

\begin{lemma}[Липшицева константа линейного слоя]
Для линейного отображения
\[
    f(\bm{x})=W\bm{x}
\]
с евклидовой нормой минимальная липшицева константа равна
\[
    \Lip(f)=\|W\|_2=\sigma_{\max}(W).
\]
Для аффинного отображения \(f(\bm{x})=W\bm{x}+\bm{b}\) результат тот
же.
\end{lemma}

\begin{proof}
Для любых \(\bm{x},\bm{y}\)
\[
    f(\bm{x})-f(\bm{y})
    =
    W(\bm{x}-\bm{y}).
\]
По определению операторной нормы
\[
    \|W(\bm{x}-\bm{y})\|_2
    \leq
    \|W\|_2\|\bm{x}-\bm{y}\|_2.
\]
Следовательно, \(\|W\|_2\) является липшицевой константой.

Покажем, что меньшая константа невозможна. По сингулярному разложению
существует единичный правый сингулярный вектор \(\bm{v}_{\max}\), для
которого
\[
    \|W\bm{v}_{\max}\|_2
    =
    \sigma_{\max}(W).
\]
Возьмем \(\bm{x}-\bm{y}=\bm{v}_{\max}\). Тогда
\[
    \frac{
        \|f(\bm{x})-f(\bm{y})\|_2
    }{
        \|\bm{x}-\bm{y}\|_2
    }
    =
    \frac{
        \|W\bm{v}_{\max}\|_2
    }{
        \|\bm{v}_{\max}\|_2
    }
    =
    \sigma_{\max}(W).
\]
Поэтому никакая константа меньше \(\sigma_{\max}(W)\) не удовлетворяет
условию Липшица.

Для аффинного слоя смещение сокращается:
\[
    (W\bm{x}+\bm{b})-(W\bm{y}+\bm{b})
    =
    W(\bm{x}-\bm{y}).
\]
Следовательно, результат не меняется.
\end{proof}

Если
\[
    F=f_L\circ f_{L-1}\circ\ldots\circ f_1,
\]
то
\[
    \Lip(F)
    \leq
    \prod_{l=1}^{L}\Lip(f_l).
\]
Поэтому контроль спектральных норм слоев ограничивает чувствительность
всей сети.

Спектральная нормализация заменяет
\[
    W
    \quad\text{на}\quad
    \overline{W}
    =
    \frac{W}{\sigma_{\max}(W)}.
\]

\begin{algobox}{Оценка спектральной нормы степенным методом}
\begin{enumerate}
    \item Инициализировать единичный вектор \(\bm{u}\).

    \item Выполнить одну или несколько итераций:
    \[
        \bm{v}
        \leftarrow
        \frac{W^\top\bm{u}}
        {\|W^\top\bm{u}\|_2},
    \]
    \[
        \bm{u}
        \leftarrow
        \frac{W\bm{v}}
        {\|W\bm{v}\|_2}.
    \]

    \item Оценить спектральную норму:
    \[
        \widehat{\sigma}
        =
        \bm{u}^\top W\bm{v}.
    \]

    \item Использовать нормированную матрицу:
    \[
        \overline{W}=\frac{W}{\widehat{\sigma}}.
    \]
\end{enumerate}
\end{algobox}

\begin{example}
Для
\[
    W=
    \begin{bmatrix}
        3&0\\
        0&1
    \end{bmatrix}
\]
спектральная норма равна \(3\). После нормализации
\[
    \overline{W}
    =
    \begin{bmatrix}
        1&0\\
        0&1/3
    \end{bmatrix},
\]
и ее спектральная норма равна \(1\).
\end{example}

% ------------------------------------------------------------
\subsection{Wasserstein GAN и градиентный штраф}
% ------------------------------------------------------------

Вместо вероятностного дискриминатора можно использовать вещественный
\emph{критик}
\[
    f_\psi:\mathcal{X}\to\R.
\]
Сигмоида на выходе критика не применяется.

Рассмотрим класс 1-липшицевых функций
\[
    \mathcal{F}_1
    =
    \left\{
        f:
        |f(x)-f(y)|
        \leq
        \|x-y\|_2
    \right\}.
\]

В WGAN используется функционал
\[
    \sup_{f\in\mathcal{F}_1}
    \left[
        \E_{x\sim P_{\mathrm{data}}}f(x)
        -
        \E_{z\sim P_z}f(G_\theta(z))
    \right].
\]

Критик максимизирует разность оценок реальных и синтетических данных.
В форме минимизации его функция потерь:
\[
\begin{split}
    L_f
    &=
    \E_{z\sim P_z}[f_\psi(G_\theta(z))]\\
    &\quad-
    \E_{x\sim P_{\mathrm{data}}}[f_\psi(x)].
\end{split}
\]

Функция потерь генератора:
\[
    L_G
    =
    -\E_{z\sim P_z}[f_\psi(G_\theta(z))].
\]

В исходном варианте WGAN ограничение Липшица приближенно обеспечивалось
ограничением весов. Это может уменьшать выразительность критика.
Более гибкий подход --- градиентный штраф.

Пусть
\[
    \widetilde{x}=G_\theta(z),
\]
\[
    \widehat{x}
    =
    \varepsilon x+(1-\varepsilon)\widetilde{x},
    \qquad
    \varepsilon\sim U[0,1].
\]

Градиентный штраф:
\[
    L_{\mathrm{GP}}
    =
    \lambda
    \E_{\widehat{x}}
    \left[
        \left(
            \|\nabla_{\widehat{x}}f_\psi(\widehat{x})\|_2-1
        \right)^2
    \right].
\]

Итоговая функция критика:
\[
\begin{split}
    L_f^{\mathrm{GP}}
    &=
    \E_{\widetilde{x}\sim P_g}[f_\psi(\widetilde{x})]
    -
    \E_{x\sim P_{\mathrm{data}}}[f_\psi(x)]\\
    &\quad+
    \lambda
    \E_{\widehat{x}}
    \left[
        \left(
            \|\nabla_{\widehat{x}}f_\psi(\widehat{x})\|_2-1
        \right)^2
    \right].
\end{split}
\]

\begin{algobox}{Алгоритм WGAN с градиентным штрафом}
\textbf{Вход:} генератор \(G_\theta\), критик \(f_\psi\), коэффициент
\(\lambda\), число обновлений критика \(n_{\mathrm{critic}}\).

\begin{enumerate}
    \item Пока генератор не достиг требуемого качества:

    \item Повторить \(n_{\mathrm{critic}}\) раз:
    \begin{enumerate}
        \item выбрать реальные объекты
        \[
            \{x_i\}_{i=1}^{B};
        \]
        \item выбрать латентные векторы
        \[
            \{z_i\}_{i=1}^{B};
        \]
        \item построить
        \[
            \widetilde{x}_i=G_\theta(z_i);
        \]
        \item выбрать
        \[
            \varepsilon_i\sim U[0,1];
        \]
        \item построить интерполяции
        \[
            \widehat{x}_i
            =
            \varepsilon_i x_i
            +(1-\varepsilon_i)\widetilde{x}_i;
        \]
        \item вычислить
        \[
        \begin{split}
            \widehat{L}_f
            &=
            \frac{1}{B}\sum_{i=1}^{B}
            \biggl[
                f_\psi(\widetilde{x}_i)-f_\psi(x_i)\\
            &\quad+
                \lambda
                \left(
                    \|\nabla_{\widehat{x}_i}
                    f_\psi(\widehat{x}_i)\|_2-1
                \right)^2
            \biggr];
        \end{split}
        \]
        \item обновить
        \[
            \psi
            \leftarrow
            \psi-\eta_f\nabla_\psi\widehat{L}_f.
        \]
    \end{enumerate}

    \item Выбрать новый batch латентных векторов и вычислить
    \[
        \widehat{L}_G
        =
        -\frac{1}{B}
        \sum_{i=1}^{B}
        f_\psi(G_\theta(z_i)).
    \]

    \item Обновить
    \[
        \theta
        \leftarrow
        \theta-\eta_G\nabla_\theta\widehat{L}_G.
    \]
\end{enumerate}
\end{algobox}

% ------------------------------------------------------------
\subsection{Hinge loss}
% ------------------------------------------------------------

Еще один распространенный критерий использует вещественный
дискриминатор \(D(x)\) без сигмоиды:
\[
\begin{split}
    L_D
    &=
    \E_{x\sim P_{\mathrm{data}}}
    [\max(0,1-D(x))]\\
    &\quad+
    \E_{z\sim P_z}
    [\max(0,1+D(G(z)))].
\end{split}
\]

Генератор минимизирует
\[
    L_G
    =
    -\E_{z\sim P_z}[D(G(z))].
\]

Объекты, уже находящиеся за требуемой границей классификации, не дают
дополнительного вклада в hinge loss дискриминатора.

% ------------------------------------------------------------
\subsection{Условные GAN}
% ------------------------------------------------------------

В условной GAN генерация зависит от дополнительной информации \(y\):
\[
    \widetilde{x}=G(z,y).
\]
Дискриминатор также получает условие:
\[
    D(x,y).
\]

Минимаксная задача:
\[
\begin{split}
    \min_G\max_D\;
    &
    \E_{(x,y)\sim P_{\mathrm{data}}}
    [\log D(x,y)]\\
    &+
    \E_{z\sim P_z,\;y\sim P_y}
    [\log(1-D(G(z,y),y))].
\end{split}
\]

Условием может быть:

\begin{itemize}
    \item метка класса;
    \item текстовое описание;
    \item другое изображение;
    \item карта сегментации;
    \item набор атрибутов;
    \item низкокачественная версия требуемого изображения.
\end{itemize}

\begin{examplebox}{Условная генерация цифр}
Пусть \(y\in\{0,\ldots,9\}\) --- требуемая цифра.

\begin{enumerate}
    \item Метка \(y\) преобразуется в embedding-вектор \(e(y)\).

    \item Генератор получает
    \[
        [z;e(y)]
    \]
    и создает изображение цифры класса \(y\).

    \item Дискриминатор получает изображение и ту же метку \(y\).

    \item Реальная пара \((x,y)\) должна получить высокую оценку.

    \item Синтетическая пара \((G(z,y),y)\) должна получить низкую
    оценку при обучении дискриминатора.
\end{enumerate}

После обучения можно фиксировать \(y\) и менять \(z\), получая разные
варианты одной цифры.
\end{examplebox}

% ------------------------------------------------------------
\subsection{Image-to-image translation}
% ------------------------------------------------------------

Пусть \(x\) --- входное изображение, а \(y\) --- требуемое выходное
изображение. Генератор строит
\[
    \widehat{y}=G(x,z).
\]

Условная состязательная функция:
\[
\begin{split}
    L_{\mathrm{cGAN}}(G,D)
    &=
    \E_{x,y}[\log D(x,y)]\\
    &\quad+
    \E_{x,z}
    [\log(1-D(x,G(x,z)))].
\end{split}
\]

Для сохранения глобального содержания добавляется реконструкционная
функция:
\[
    L_{L_1}(G)
    =
    \E_{x,y,z}
    \left[
        \|y-G(x,z)\|_1
    \right].
\]

Итоговая задача:
\[
    \min_G\max_D
    \left[
        L_{\mathrm{cGAN}}(G,D)
        +
        \lambda L_{L_1}(G)
    \right].
\]

Состязательная часть отвечает преимущественно за реалистичность, а
\(L_1\)-часть --- за соответствие целевому изображению.

Для генератора часто используется U-Net. Skip connections передают
детальную пространственную информацию из encoder в decoder.

PatchGAN-дискриминатор выдает не один скаляр, а сетку оценок:
\[
    D(x,y)\in\R^{h\times w}.
\]
Каждая оценка относится к локальному фрагменту изображения. Итоговая
функция потерь усредняется по всем фрагментам.

% ------------------------------------------------------------
\subsection{Оценка качества GAN}
% ------------------------------------------------------------

Функции потерь генератора и дискриминатора не всегда непосредственно
соответствуют визуальному качеству. Поэтому используются отдельные
метрики.

\subsubsection{Визуальная оценка}

Проверяются:

\begin{itemize}
    \item реалистичность отдельных объектов;
    \item разнообразие batch;
    \item наличие повторяющихся объектов;
    \item плавность интерполяций в латентном пространстве;
    \item ближайшие соседи среди обучающих данных;
    \item качество генерации для фиксированных латентных векторов на
    разных шагах обучения.
\end{itemize}

Латентная интерполяция:
\[
    z(t)=(1-t)z_1+tz_2,
    \qquad t\in[0,1].
\]
Плавное изменение \(G(z(t))\) свидетельствует о регулярной структуре
изученного латентного пространства.

\subsubsection{Inception Score}

Пусть предварительно обученный классификатор задает \(p(y\mid x)\).
Определим
\[
    p(y)=\E_x[p(y\mid x)].
\]
Тогда
\[
    \operatorname{IS}
    =
    \exp\left(
        \E_x
        \KL\bigl(
            p(y\mid x)\|p(y)
        \bigr)
    \right).
\]

Высокий IS требует:

\begin{itemize}
    \item малого условного энтропийного разброса \(p(y\mid x)\), то
    есть уверенно классифицируемых изображений;
    \item разнообразного маргинального распределения \(p(y)\).
\end{itemize}

Недостатки: метрика не сравнивает синтетические данные напрямую с
реальными и зависит от используемого классификатора.

\subsubsection{Fréchet Inception Distance}

Пусть признаки реальных и синтетических изображений приближенно
моделируются гауссовскими распределениями
\[
    \mathcal{N}(\mu_r,\Sigma_r),
    \qquad
    \mathcal{N}(\mu_g,\Sigma_g).
\]

Тогда
\[
\begin{split}
    \operatorname{FID}
    &=
    \|\mu_r-\mu_g\|_2^2\\
    &\quad+
    \tr\left(
        \Sigma_r+\Sigma_g
        -
        2
        \left(
            \Sigma_r^{1/2}
            \Sigma_g
            \Sigma_r^{1/2}
        \right)^{1/2}
    \right).
\end{split}
\]

Меньшее значение считается лучшим. FID сравнивает реальные и
синтетические данные, но:

\begin{itemize}
    \item зависит от пространства признаков;
    \item имеет статистическую погрешность при малом числе примеров;
    \item сравнивает только первые два момента признаков;
    \item должен вычисляться с одинаковой предобработкой изображений.
\end{itemize}

\subsubsection{Precision и recall для генеративных моделей}

В неформальной интерпретации:

\begin{itemize}
    \item \textbf{precision} измеряет, насколько синтетические объекты
    похожи на допустимые реальные объекты;
    \item \textbf{recall} измеряет, насколько полно генератор
    покрывает разнообразие реального распределения.
\end{itemize}

Высокая реалистичность при низком разнообразии может давать высокий
precision, но низкий recall.

% ------------------------------------------------------------
\subsection{Диагностика обучения GAN}
% ------------------------------------------------------------

\begin{center}
\begin{tabularx}{\textwidth}{
    >{\bfseries}p{3.3cm}
    X
    X
}
\toprule
Симптом & Возможная причина & Возможные действия\\
\midrule

Все образцы похожи
&
Коллапс мод.
&
Усилить регуляризацию дискриминатора, изменить функцию потерь,
проверить баланс обновлений, использовать minibatch-признаки или
градиентный штраф.\\

Дискриминатор почти идеален
&
Он обучается существенно быстрее генератора.
&
Уменьшить его скорость обучения или число обновлений, усилить
аугментацию, применить спектральную нормализацию.\\

Оба участника не улучшаются
&
Слишком слабый дискриминатор, неверная архитектура или ошибка в
градиентах.
&
Проверить отделение графов вычислений, знаки функций потерь,
нормализацию входов и выходную нелинейность.\\

Появляются NaN
&
Взрыв градиента, численно неустойчивые логарифмы, слишком большая
скорость обучения.
&
Использовать функции потерь по логитам, уменьшить скорость обучения,
добавить gradient clipping или спектральную нормализацию.\\

Шахматный узор
&
Неравномерное перекрытие транспонированных сверток.
&
Использовать upsampling с последующей сверткой.\\

Хорошая метрика, но много копий
&
Переобучение или недостаточная проверка разнообразия.
&
Сравнить с ближайшими обучающими объектами, вычислить recall и
проверить интерполяции.\\

Резкие скачки качества
&
Нестабильная динамика игры.
&
Использовать EMA генератора, меньшие скорости обучения, WGAN-GP,
спектральную нормализацию или разные скорости участников.\\
\bottomrule
\end{tabularx}
\end{center}

% ------------------------------------------------------------
\subsection{Сравнение основных вариантов GAN}
% ------------------------------------------------------------

\begin{center}
\begin{tabularx}{\textwidth}{
    >{\bfseries}p{2.7cm}
    X
    X
}
\toprule
Модель & Основная идея & Особенности\\
\midrule

Классическая GAN
&
Бинарная классификация реальных и синтетических данных.
&
Минимаксная или ненасыщающаяся функция потерь.\\

DCGAN
&
Сверточные генератор и дискриминатор.
&
Постепенное увеличение и уменьшение пространственного разрешения.\\

Conditional GAN
&
Генерация зависит от условия \(y\).
&
Условие передается генератору и дискриминатору.\\

Pix2Pix
&
Условное преобразование изображения в изображение.
&
Состязательная функция плюс \(L_1\), U-Net и PatchGAN.\\

WGAN
&
Вещественный критик и липшицево ограничение.
&
Нет сигмоиды; критик ранжирует реальные и синтетические объекты.\\

WGAN-GP
&
Мягкое обеспечение липшицевости.
&
Штраф за отклонение нормы градиента от единицы.\\

SN-GAN
&
Спектральная нормализация дискриминатора.
&
Контроль операторных норм слоев.\\
\bottomrule
\end{tabularx}
\end{center}

% ============================================================
\section{Связь CNN и GAN}
% ============================================================

CNN и GAN описывают разные уровни построения модели:

\begin{itemize}
    \item CNN --- семейство архитектурных преобразований;
    \item GAN --- способ обучения генеративной модели как игры двух
    сетей.
\end{itemize}

GAN не обязана быть сверточной. Генератор и дискриминатор могут быть
полносвязными, рекуррентными или другими моделями. Однако для
изображений обычно используются сверточные архитектуры.

\begin{center}
\begin{tabularx}{\textwidth}{
    >{\bfseries}p{3.4cm}
    X
    X
}
\toprule
Свойство & Классификационная CNN & GAN\\
\midrule

Основная цель
&
Предсказать метку или структурированный выход.
&
Приблизить распределение данных и генерировать новые объекты.\\

Компоненты
&
Одна сеть.
&
Генератор и дискриминатор или критик.\\

Тип обучения
&
Обычно минимизация одной функции потерь.
&
Состязательная игра с чередующимися обновлениями.\\

Вход
&
Изображение или другой сигнал.
&
Генератор получает шум и, возможно, условие; дискриминатор получает
реальный или синтетический объект.\\

Выход
&
Классы, координаты, маски или признаки.
&
Синтетические данные и оценка их реалистичности.\\

Типичные операции
&
Свертки, pooling, нормализация, residual-блоки.
&
Свертки с уменьшением разрешения в дискриминаторе и увеличение
разрешения в генераторе.\\
\bottomrule
\end{tabularx}
\end{center}

% ============================================================
\section{Краткий план воспроизведения материала}
% ============================================================

\subsection*{Сверточная сеть}

Последовательность основных понятий:

\begin{enumerate}
    \item входной тензор \(C\times H\times W\);
    \item локальное ядро и разделение параметров;
    \item формула многоканальной свертки;
    \item padding, stride, dilation;
    \item размер выхода;
    \item число параметров;
    \item эквивариантность относительно сдвига;
    \item рецептивное поле;
    \item нелинейность, pooling и нормализация;
    \item специальные свертки;
    \item residual connections;
    \item функция потерь и обратное распространение;
    \item обучение mini-batch-методами;
    \item регуляризация и аугментация.
\end{enumerate}

\subsection*{Генеративная состязательная сеть}

Последовательность основных понятий:

\begin{enumerate}
    \item латентная переменная \(z\sim P_z\);
    \item генератор \(G(z)\) и распределение \(P_g\);
    \item дискриминатор \(D(x)\);
    \item минимаксная функция;
    \item оптимальный дискриминатор
    \[
        D^*(x)
        =
        \frac{p_{\mathrm{data}}(x)}
        {p_{\mathrm{data}}(x)+p_g(x)};
    \]
    \item связь оптимального критерия с дивергенцией
    Дженсена--Шеннона;
    \item глобальный оптимум \(P_g=P_{\mathrm{data}}\);
    \item ненасыщающаяся функция генератора;
    \item чередующиеся обновления;
    \item mode collapse и нестабильность;
    \item сверточная архитектура генератора и дискриминатора;
    \item спектральная нормализация;
    \item WGAN-GP;
    \item условные GAN;
    \item FID, IS, precision и recall.
\end{enumerate}

% ============================================================
\section{Итоговые положения}
% ============================================================

\begin{enumerate}
    \item Сверточный слой использует локальность и разделение
    параметров, благодаря чему число параметров не зависит от
    пространственного размера входа.

    \item Свертка с единичным шагом без граничных эффектов
    эквивариантна относительно сдвигов.

    \item Глубина сети и операции уменьшения разрешения увеличивают
    рецептивное поле.

    \item Остаточные связи создают прямые пути распространения
    признаков и градиентов.

    \item GAN задает неявную генеративную модель через преобразование
    латентной случайной переменной.

    \item Для фиксированного генератора оптимальный дискриминатор
    определяется отношением плотностей реального и генерируемого
    распределений.

    \item При неограниченной мощности моделей глобальный оптимум
    классической GAN достигается при
    \[
        P_g=P_{\mathrm{data}},
        \qquad
        D^*(x)=\frac{1}{2}.
    \]

    \item Теоретический глобальный оптимум не гарантирует, что
    градиентный алгоритм достигнет его при конечной выборке, конечной
    мощности сетей и неидеальной оптимизации.

    \item Основные практические проблемы GAN --- нестабильность,
    коллапс мод, дисбаланс генератора и дискриминатора, артефакты и
    сложность объективной оценки качества.

    \item Для стабилизации применяются ненасыщающиеся функции потерь,
    спектральная нормализация, градиентные штрафы, разные скорости
    обучения и экспоненциальное усреднение параметров генератора.
\end{enumerate}

\end{document}
